\DocumentMetadata{
  pdfversion=2.0,
  pdfstandard=ua-2,
  lang=en-US,
  tagging=on,
  tagging-setup={math/setup=mathml-SE,tabsorder=structure}
}
\documentclass[11pt,letterpaper]{article}
\usepackage[margin=0.68in,includefoot]{geometry}
\usepackage{newcomputermodern}
\usepackage{amsmath,amscd,mathtools}
\providecommand{\square}{\mdlgwhtsquare}
\usepackage[hidelinks]{hyperref}
\hypersetup{
  pdftitle={Numerical Analysis PhD Exam, January 2017},
  pdfauthor={University of Florida Department of Mathematics},
  pdfsubject={Graduate mathematics examination},
  pdfdisplaydoctitle=true,
  bookmarksopen=true
}
\setcounter{secnumdepth}{0}
\setlength{\parindent}{0pt}
\setlength{\parskip}{0.28em}
\setlength{\emergencystretch}{3em}
\setlength{\leftmargini}{2.15em}
\setlength{\leftmarginii}{2.65em}
\setlength{\leftmarginiii}{3em}
\renewcommand{\labelenumi}{\textbf{\theenumi.}}
\pagestyle{empty}
\raggedbottom
\allowdisplaybreaks
\newenvironment{examproblems}[1]{%
  \begin{enumerate}%
  \setcounter{enumi}{#1}%
  \setlength{\topsep}{0.35em}%
  \setlength{\partopsep}{0pt}%
  \setlength{\itemsep}{0.48em}%
  \setlength{\parsep}{0.18em}%
}{\end{enumerate}}
\newenvironment{examsubparts}{%
  \begin{enumerate}%
  \setlength{\topsep}{0.18em}%
  \setlength{\partopsep}{0pt}%
  \setlength{\itemsep}{0.16em}%
  \setlength{\parsep}{0.1em}%
}{\end{enumerate}}
\newenvironment{examinstructions}{%
  \par\begingroup\itshape
}{%
  \par\endgroup\vspace{0.18em}
}
\makeatletter
\renewcommand\section{\@startsection{section}{1}{0pt}%
  {-1.25ex plus -.25ex minus -.1ex}%
  {0.55ex plus .1ex}%
  {\normalfont\large\bfseries}}
\renewcommand{\@maketitle}{%
  \newpage\null\vskip -1.2em%
  {\footnotesize\bfseries
    \href{https://www.ufl.edu/}{University of Florida}, \href{https://math.ufl.edu/}{Department of Mathematics}\par}%
  \vskip 0.18em%
  \tagstructbegin{tag=Title}%
    {\LARGE\bfseries\href{https://gma.math.ufl.edu/exams/numerical-analysis-phd/}{Numerical Analysis PhD Exam}, January 2017\par}%
  \tagstructend%
  \vskip 0.35em%
  \tagmcbegin{artifact}%
    \hrule height 1pt%
  \tagmcend%
  \vskip 0.55em%
}
\makeatother
\title{Numerical Analysis PhD Exam, January 2017}
\author{}
\date{}
\begin{document}
\maketitle
\thispagestyle{empty}
\begin{examinstructions}
You must show all your work as neatly and clearly as possible and indicate the final answer clearly. You may not use a calculator.
\end{examinstructions}
\begin{examproblems}{0}
\item
This problem has two parts.
\begin{examsubparts}
\item[\textnormal{(a)}]
Show that the equation \(x=\frac{1}{2}\cos(x)\) has a solution~\(\alpha\).
\item[\textnormal{(b)}]
Find an interval \([a,b]\) containing~\(\alpha\) and such that for every \(x_0\in[a,b]\), the iterative sequence 
\[
x_{n+1}=\frac{1}{2}\cos x_n
\]
 will converge to~\(\alpha\). Justify your answer.
\end{examsubparts}
\item
For the basic Lagrange polynomials 
\[
L_i(x)=\prod_{j\ne i}\frac{x-x_j}{x_i-x_j}\qquad\text{for}\qquad i=0,\ldots,n,
\]
 show that 
\[
\sum_i L_i(x)=1\qquad\text{for all }x.
\]
\item
Consider a quadrature rule of the form for \(0<\alpha<1\): 
\[
\int_0^1 x^\alpha f(x)\,dx\approx A\int_0^1 f(x)\,dx+B\int_0^1 xf(x)\,dx.
\]
\begin{examsubparts}
\item[\textnormal{(a)}]
Determine the constants \(A,B\) so that the quadrature formula has maximum degree of exactness.
\item[\textnormal{(b)}]
What is the degree of exactness of the formula in part (a)?
\end{examsubparts}
\item
Let~\(f\) be an arbitrary (continuous) function on \([0,1]\) satisfying 
\[
f(x)+f(1-x)\equiv 1\qquad\text{for}\qquad 0\le x\le 1.
\]
\begin{examsubparts}
\item[\textnormal{(a)}]
Show that \(\int_0^1 f(x)\,dx=\frac{1}{2}\).
\item[\textnormal{(b)}]
Show that the composite trapezoidal rule for computing \(\int_0^1 f(x)\,dx\) is exact.
\end{examsubparts}
\item
Assume that you are solving the initial value problem 
\[
\begin{aligned}
y'&=f(t,y), && a\le t\le b,\\
y(a)&=\alpha.
\end{aligned}
\]
 The formula for the global error of the numerical solutions for the ODE problem above obtained via Euler’s method is (\(M=\lVert Y''\rVert_\infty\)): 
\[
\lvert Y(t_i)-w_i\rvert<\frac{hM}{2L}\left[e^{L(b-a)}-1\right].
\]
 Compute the values of~\(L\) and \(M=\lVert Y''\rVert_\infty\) necessary to apply the global error formula above to the specific ODE problem 
\[
\begin{aligned}
y'&=\sin(t+2y)+e^t, && 0\le t\le 1,\\
y(0)&=0.
\end{aligned}
\]
\end{examproblems}
\end{document}
